% !TeX program = xelatex
\documentclass[11pt,a4paper]{article}
\usepackage[margin=23mm]{geometry}
\usepackage{fontspec}
\setmainfont{Latin Modern Roman}
\usepackage{amsmath,amssymb,mathtools}
\usepackage{xcolor,enumitem,booktabs,longtable,array,xurl,hyperref}
\definecolor{muted}{HTML}{53636C}
\hypersetup{colorlinks=true,urlcolor=blue,linkcolor=blue}
\setlength{\parindent}{0pt}
\setlength{\parskip}{5pt}
\newcommand{\R}{\mathbb R}
\newcommand{\E}{\mathbb E}
\newcommand{\N}{\mathbb N}
\newcommand{\ind}{\mathbf 1}
\newcommand{\Var}{\operatorname{Var}}
\newcommand{\Cov}{\operatorname{Cov}}
\newcommand{\supp}{\operatorname{supp}}
\DeclareMathOperator*{\argmin}{arg\,min}
\DeclareMathOperator*{\argmax}{arg\,max}
\newcommand{\entrymeta}[3]{{\sffamily\small\color{muted}\textbf{\texttt{#1}}\quad|\quad #2\par #3\par}}
\newcommand{\Problemref}[1]{\texttt{#1}}
\newcommand{\JELref}[2]{\texttt{#2} (JEL \texttt{#1})}
\begin{document}
\section*{Minimax estimation of counterfactual outcome densities}
\textbf{Problem ID:} \texttt{C14-111}\quad \textbf{JEL:} \texttt{C14}\par
% BEGIN ECONOMICS STATEMENT
\label{problem:OP-0529}
\label{beyond:CI-50}
\entrymeta{CI-50}{Econometric theory and causal inference}{Research agenda formalization}
\subsection*{Definitions and assumptions}
Observe $n$ independent and identically distributed copies $O_1,\ldots,O_n$ of $O=(X,W,Y)$ with common observable law $P$, where $X\in[0,1]^d$, $W\in\{0,1\}$, and continuous $Y\in[0,1]$. The data law is $P^{\otimes n}$; potential outcomes are not observed in addition to these data. Assume consistency $Y=Y(W)$, conditional exchangeability $(Y(0),Y(1))\perp W\mid X$, and fixed overlap $\epsilon\leq e(x)\leq1-\epsilon$, where $e(x)=\Pr_P(W=1\mid X=x)$ and $0<\epsilon<1/2$. Let $p$ be the Lebesgue density of $X$ on $[0,1]^d$, normalized by $\int_{[0,1]^d}p(x)\,dx=1$. Fix $w\in\{0,1\}$ and let $f_w(y\mid x)$ be the Lebesgue density of the conditional law of $Y$ given $W=w,X=x$. The functions $p,e,f_w$ are unknown. Fix $d\in\N$ with $d\geq1$, positive $s_f,s_e,s_p,L,c,C$, and $\epsilon$ independently of $n$, and require
\[
 f_w\in\mathcal H_{d+1}^{s_f}(L),\quad
 e\in\mathcal H_d^{s_e}(L),\quad
 p\in\mathcal H_d^{s_p}(L),\quad
 c\leq p\leq C,\quad 0\leq f_w\leq C.
\]
Use the continuous versions supplied by these H\"older classes on their closed cubes, and impose the displayed bounds and overlap on those versions. Each $f_w(\cdot\mid x)$ integrates to one for every $x\in[0,1]^d$. The other treatment arm is unrestricted apart from the stated support. Let $\mathcal P$ be the class of observable laws satisfying these restrictions and admitting potential outcomes with the stated consistency and exchangeability properties; no potential-outcome joint law beyond these properties is specified. No nuisance function is supplied to the estimator.

Keep the declared constant choices rather than assuming strict slack around uniform densities. Empty classes are to be identified as such and do not define an estimation problem here. For nonempty classes, distinguish cases in which the target is known throughout the class and the risks are zero from cases with a positive minimax rate; a positive two-sided rate is requested only in the latter cases. Constants in rate comparisons may depend on $d,s_f,s_e,s_p,L,c,C,\epsilon$ and, for the supremum loss, $\eta$, but not on $n$.

The counterfactual density is identified as
$f_{Y(w)}(y)=\int_{[0,1]^d} f_w(y\mid x)p(x)\,dx$.
\subsection*{Formal question}
For each nonempty $\mathcal P$, determine matching minimax rates for estimating $f_{Y(w)}$ under integrated squared error,
\[
 R_{n,2}=\inf_{\widehat f_n}\sup_{P\in\mathcal P}
 \E_{P^{\otimes n}}\int_0^1\{\widehat f_n(y)-f_{Y(w)}(y)\}^2\,dy.
\]
Separately, for fixed $\eta>0$ with nonempty $I_\eta=[\eta,1-\eta]$, determine the squared supremum risk
\[
 R_{n,\infty}(\eta)=\inf_{\widehat f_n}\sup_{P\in\mathcal P}
 \E_{P^{\otimes n}}\left[\sup_{y\in I_\eta}
 |\widehat f_n(y)-f_{Y(w)}(y)|^2\right].
\]
The infima range over measurable real-valued function estimators based only on $O_1,\ldots,O_n$ and optional independent randomization, with measurable displayed losses; estimators need not themselves be normalized densities. For $0<\eta<1/2$ the second loss is uniform on a positive-length interior interval; for $\eta=1/2$ it is squared error at $y=1/2$. For $\eta>1/2$ the stated interval is empty, and no supremum-loss convention or estimation question on that empty domain is specified.

Describe when marginalization over $X$ permits a rate governed primarily by the one-dimensional outcome, and when propensity or covariate-law complexity changes that rate. Lower-bound constructions must be valid submodels of the observable experiment $\{P^{\otimes n}:P\in\mathcal P\}$ and match each claimed regime. Ordinary density-estimation submodels may be used where they suffice; any claimed additional nuisance-driven difficulty requires a lower bound establishing that difficulty in this experiment.
\subsection*{Scope and status}
This question concerns treatment-specific marginal densities, not the generally unidentified density of $Y(1)-Y(0)$, which also requires their joint law. Rates for quantiles or the number of modes require additional regularity and are not consequences of density identification alone. The model is an explicit agenda formulation.
\subsection*{References for the source of the problem}
Carlos Cinelli, Avi Feller, Guido Imbens, Edward Kennedy, Sara Magliacane, and Jose Zubizarreta, \emph{Challenges in Statistics: A Dozen Challenges in Causality and Causal Inference}, arXiv:2508.17099v1, 23 August 2025, Section 3.5.3, \emph{New parameters}.
\url{https://arxiv.org/html/2508.17099v1#S3.SS5.SSS3}

\subsection*{Common conventions: Source mathematical conventions}

\textbf{Well-defined objects.}
All probability spaces and measurable variables are specified or inherited
from a local statistical model. Expectations used as finite targets require
the stated integrability. A decision rule or estimator is measurable with
respect to its available information; randomized rules may use an auxiliary
independent random variable. Norms without a subscript are Euclidean in
finite-dimensional spaces. An optimizer is requested only when attainment
is assured; otherwise the question concerns the optimal value and
$\varepsilon$-optimal admissible rules for every $\varepsilon>0$.

\textbf{Identification.}
A structural model $m\in\mathcal M$ implies an observable law
$P_m^{\rm obs}$ and target $\theta(m)$. The target is point identified on
$\mathcal M$ if
\[
 P_{m_1}^{\rm obs}=P_{m_2}^{\rm obs}
 \quad\Longrightarrow\quad\theta(m_1)=\theta(m_2)
 \qquad(m_1,m_2\in\mathcal M).
\]
When this fails, the sharp identified set is
\[
 \Theta(P;\mathcal M)=\{\theta(m):m\in\mathcal M,
                                  P_m^{\rm obs}=P\}.
\]
Specifying an intervention or estimand does not establish identification.
Positivity, exclusion, causal sufficiency, measurement restrictions, and
transport assumptions are substantive local inputs, never automatic
consequences of notation.

\textbf{Minimax risk and nontrivial inference.}
For a statistical experiment with data law $P^{(n)}$, target $\theta(P)$,
and nonnegative loss $\ell$, define
\[
 \mathcal R_n(\mathcal P,\ell)
   =\inf_{\widehat\theta_n}\sup_{P\in\mathcal P}
      \E_{P^{(n)}}\ell(\widehat\theta_n,\theta(P)).
\]
The infimum is over measurable estimators. A sharp rate is a sequence
$r_n$ for which $0<c\leq C<\infty$, independent of $n$, satisfies
$c r_n\leq\mathcal R_n\leq C r_n$ for all sufficiently large $n$.
Squared-error parametric risk is of order $n^{-1}$; the corresponding
root-mean-squared error is of order $n^{-1/2}$. These different conventions
are distinguished locally. A uniform asymptotic confidence region $C_n$
must satisfy
\[
 \liminf_{n\to\infty}\inf_{P\in\mathcal P}
     \Pr_{P^{(n)}}\{\theta(P)\in C_n\}\geq 1-\alpha,
 \qquad 0<\alpha<1,
\]
together with an informative diameter or expected-width requirement.
Coverage alone permits the vacuous whole parameter space. Testing questions
similarly impose both size control and power against a declared separated
alternative class.

\textbf{Binary causal baseline.}
When an entry explicitly invokes this baseline, one observes independent
copies of $O=(X,W,Y)$, with $W\in\{0,1\}$ and potential outcomes $Y(0),Y(1)$.
Assume consistency $Y=Y(W)$, no interference, conditional exchangeability
$(Y(0),Y(1))\perp W\mid X$, and overlap
$\epsilon\leq e(X)\leq1-\epsilon$ almost surely for a fixed
$0<\epsilon<1/2$, where
\[
 e(x)=\Pr(W=1\mid X=x),\qquad
 m_w(x)=\E[Y\mid W=w,X=x],\qquad
 \tau(x)=m_1(x)-m_0(x).
\]
These assumptions identify conditional mean effects almost everywhere
under the covariate law. Pointwise or uniform effects require appropriate
regularity and a specified support. Continuous treatment, longitudinal
data, adaptive experiments, and network interference require their own
assumptions in place of this baseline. Bounded outcomes, densities, and
smoothness are additional assumptions only where stated.

\textbf{Function classes.}
On $[0,1]^d$, $\mathcal H_d^s(L)$ is the isotropic Hölder ball of radius
$L>0$: put $k=\lceil s\rceil-1$ and $\gamma=s-k\in(0,1]$, bound derivatives
through order $k$, and require the order-$k$ derivatives to be
$\gamma$-Hölder with constant at most $L$. Thus integer orders use a
Lipschitz final derivative convention. The class
$\operatorname{BV}_d(L)$ consists of functions $f\in L^1((0,1)^d)$ whose
distributional derivative is a finite vector measure and for which
$\|f\|_{L^1}+|Df|((0,1)^d)\leq L$.
An $L^\infty$ bound or a particular pointwise version is imposed separately
when needed. Sparse and neural-network classes require a declared
dictionary or architecture and coefficient bounds.

An anisotropic H\"older class with declared block smoothness indices means
coordinatewise H\"older regularity in each block, uniformly in the other
blocks, with all corresponding derivative and H\"older bounds at most
the stated radius. Derivative orders and the integer-order convention in
each block are those just given. Mixed-derivative bounds are additional
restrictions only when explicitly stated.

\textbf{Decision and computational questions.}
Policy regret compares admissible feasible policies under the same law and
constraints. In computational entries, finite data and thresholds have
rational encodings; running time is measured in their total bit length.
Real-valued equilibrium search uses the explicitly declared algebraic or
FIXP encoding. An approximation ratio for maximization is written
$\operatorname{OPT}/\operatorname{ALG}$ and for minimization as
$\operatorname{ALG}/\operatorname{OPT}$, with zero optima and infeasible
instances handled locally. Historical approximation bounds are not
automatically current bounds.
% END ECONOMICS STATEMENT
\end{document}
